\documentclass[a4paper,10pt]{report}\input{../0.Préambule/Préambule_cours_prof}\begin{document}

\definecolor{titlecolor}{rgb}{0.12, 0.56, 1.0}
\newpage
\setcounter{chapter}{6}
\chapter{Expressions littérales 1}

\vspace{-1mm}
\section{Chercher un problème}

\vspace{-3mm}
\begin{ex}

\noindent \begin{minipage}{6.5cm}

Voici un programme de calcul :

\medskip
\arrayrulecolor{black}
\renewcommand{\arraystretch}{0.1}
\begin{tabular}{|p{6.5cm}|}
\hline
\vspace{2mm} \edupythoninline{Choisis un nombre} \\ 
\vspace{2mm} \edupythoninline{Ajoute 3} \\ 
\vspace{2mm} \edupythoninline{Multiplie le resultat par 2} \\ 
\vspace{2mm} \edupythoninline{Ajoute 4} \\ 
\vspace{2mm} \edupythoninline{Retranche deux fois le nombre de depart} \\
\vspace{2mm} \\
\hline
\end{tabular}
\end{minipage}
\hspace{3mm}
\begin{minipage}{9.8cm}
\begin{enumerate}
\item Choisis un nombre, effectue ce programme de calcul.
Recommence avec d'autres nombres. Que constates-tu ?
\item Écris une expression numérique traduisant ce programme pour le nombre que tu as choisi.
\item On appelle N le nombre que tu as choisi. Dans l'expression trouvé au 2), remplace le nombre que tu as choisi par la lettre N.
Tu obtiens ce que l'on appelle \og une expression littérale \fg{}.
\end{enumerate}
\end{minipage}

\medskip
\noindent\grandscarreaux{\linewidth}{3.2cm}
\end{ex}


\begin{meth}
Pour chercher un problème de ce type :
\begin{enumerate}
\item Je vérifie que je comprends :

\vspace{-7.5mm}
\begin{enumerate}
\begin{tabularx}{\linewidth}{p{3.6cm}>{\arraybackslash}p{2.3cm}>{\arraybackslash}X}
&
\item[•] l'énoncé.&
\item[•] ce que l'on me demande, le type de réponse attendue.\\
\end{tabularx}
\end{enumerate}
\item Si besoin j'explicite mon problème en le représentant et/ou en le modélise
\begin{enumerate}
\item[•] Je fais les schémas, je traduis le problème en langage mathématiques.
\end{enumerate}
\item Je cherche ce qu'il se passe dans les cas simples,
\begin{enumerate}
\item[•] Je fais les calculs pour les premiers cas ....
\end{enumerate}
\item Quand je pense avoir une idée, je fais une proposition de solution.
\item Je vérifie si ma proposition fonctionne sur les cas simples.
\begin{enumerate}
\item[•] Sinon je réfléchis à une modification pour faire une nouvelle proposition. (retour au 1.)
\end{enumerate}
\item Quand je pense avoir trouvé, je rédige ma réponse en expliquant comment j'ai fait pour
trouver et pourquoi cela semble correct.
\end{enumerate}
\end{meth}

\vspace{-3mm}
\section{Produire une expression littérale}

\vspace{-2mm}
\subsection{Définition d'une expression littérale}

\begin{dft}
Une expression littérale est une expression contenant des lettres qui remplacent
des nombres.

\noindent Elle sert, par exemple, à :

\vspace{-2mm}
\begin{enumerate}
\begin{tabularx}{\linewidth}{*{4}{>{\arraybackslash}X}}
\item[•] Établir une formule &
\item[•] Traduire l'énoncé &
\item[•] Trouver un nombre &
\item[•] Prouver un résultat. \\
\end{tabularx}
\end{enumerate}
\end{dft}

\begin{ex}
\begin{enumerate}
\item La formule du périmètre d'un cercle est $2 \times r \times \pi$ ou $d \times \pi$ .
Dans ces expressions numériques $r$ est le rayon du cercle $d$ est le diamètre et $\pi$ est un
nombre qui ne change pas est qui vaut environ $3,14$.
\item On obtient la vitesse avec la formule suivante : $V = \dfrac{d}{t}$ ou $V$ est la vitesse, $d$ la distance et $t$ le temps.
\end{enumerate}
\end{ex}

\vspace{-3mm}
\subsection{Convention d'écriture}

\vspace{-2mm}
\noindent \begin{minipage}{11cm}
\begin{dft}
$5 \times 5 = 5^2$ lire \og $5$ au carré \fg{} ; \quad de même $7 \times 7\times 7 = 7^3$ lire \og $7$ au cube \fg{}

\medskip
\noindent On peut faire de même avec les lettres :

\noindent $x \times x = x^2$ lire \og $x$ au carré \fg{} ; \quad de même $x \times x \times x = x^3$ lire \og $x$ au cube \fg{}
\end{dft}

\begin{regle}
Pour simplifier l'écriture d'une expression littérale, on peut supprimer le symbole \og $\times$ \fg{} devant une lettre ou une parenthèse.
\end{regle}

\begin{rmq}
On ne peut pas supprimer le signe \og $\times$ \fg{} entre deux nombres écrits à l'aide de chiffres.
En effet, $4 \times 5$ est différent de $45$.

\medskip
\noindent On a les résultats suivants : \quad $1 \times x = x$; \quad $0 \times x = 0$; \quad $x \times y = y \times x$
\end{rmq}

\end{minipage}
\hspace{3mm}
\begin{minipage}{7cm}
\begin{ex}
On peut supprimer le signe $\times$ entre :
\begin{enumerate}
\item[•]Deux lettres, $$d \times \pi = d \pi$$
\item[•]Entre un nombre et une lettre $$2 \times r = 2 r$$
\item[•]Devant une parenthèse. $$5 \times (2 + x) = 5 (2 + x)$$
$$c \times (x + 4) = c(x + 4)$$
\end{enumerate}
\end{ex}
\end{minipage}

\subsection{Vocabulaire}

\vspace{-2mm}
\begin{dft}
Écrire un résultat « en fonction de $x$ », c'est écrire une expression littérale
contenant la lettre $x$.
\end{dft}

\begin{ex} 
\vspace{-5.5mm} \hspace{2.8cm} J'ai choisi un nombre $x$. Je l'ai multiplié par $3$, puis j'ai ajouté $7$.
Écris le résultat en fonction de $x$.

\medskip
\noindent\grandscarreaux{\linewidth}{1.6cm}
\end{ex}

\vspace{-3mm}
\section{Calculer la valeur d'une expression littérale}

\vspace{-2mm}
\begin{dft}
Calculer la valeur d'une expression littérale, c'est attribuer un nombre à chaque lettre de l'expression afin d'effectuer le calcul
\end{dft}

\begin{rmq}

\vspace{-4.5mm} \hspace{3cm} Une même lettre désigne toujours un même nombre dans une expression littérale donnée.
\end{rmq}

\vspace{-2mm}
\begin{meth}
Une golfeuse frappe une balle de golf.
La hauteur $h$, en mètre, à laquelle se trouve cette balle, $t$ secondes après son départ, est
donnée par la formule $h = 30 \times t - 4,9 \times t^2$

\medskip
\noindent Cette formule est valable jusqu'au moment où la balle retombe au sol.

\medskip
\noindent Calcule la hauteur à laquelle se trouve la balle 5 secondes après le lancer.

\medskip
\noindent\grandscarreaux{\linewidth}{4cm}
\end{meth}
\end{document}